\clearpage\pagestyle{otMain}\fontsize{11}{15}\selectfont\corrigesChapitre{13}{Thalès, réciproque et contraposée}
\begin{blocLivre}{cor-13-1}\begin{corrige}{1}{Écrire les rapports}\demarche{Garder la même correspondance des sommets.}
Les alignements et le parallélisme permettent d’appliquer le théorème de Thalès : $\res{AM/AB=AN/AC=MN/BC}$. Les côtés du petit triangle sont tous au numérateur.
\label{fig-cor-13-1}\figureGeo[corrige]{thales-01}
\end{corrige}\end{blocLivre}
\begin{blocLivre}{cor-13-2}\begin{corrige}{2}{Une longueur réduite}\demarche{Multiplier la grande longueur par le rapport de réduction.}
Thalès donne $AN/10=3/5$. Ainsi $\res{AN=6\text{ cm}}$. Cette longueur est inférieure à AC, conformément à l’emboîtement.
\label{fig-cor-13-2}\figureGeo[corrige]{thales-02}
\end{corrige}\end{blocLivre}
\begin{blocLivre}{cor-13-3}\begin{corrige}{3}{La petite base}\demarche{Choisir le rapport portant l’inconnue.}
$MN/15=4/10$, donc $\res{MN=6\text{ cm}}$. Le rapport des longueurs est 0,4.
\label{fig-cor-13-3}\figureGeo[corrige]{thales-03}
\end{corrige}\end{blocLivre}
\begin{blocLivre}{cor-13-4}\begin{corrige}{4}{Retrouver le grand côté}\demarche{Une longueur agrandie se trouve en divisant par le rapport.}
$6/AB=4/10$. Le produit en croix donne $4AB=60$, donc $\res{AB=15\text{ cm}}$.
\label{fig-cor-13-4}\figureGeo[corrige]{thales-04}
\end{corrige}\end{blocLivre}
\begin{blocLivre}{cor-13-5}\begin{corrige}{5}{De la partie au tout}\demarche{Reconstituer le côté entier avant Thalès.}
Comme M est entre A et B, AB = 3 + 2 = 5 cm. Puis $AN=12\times3/5=\res{7{,}2\text{ cm}}$. Employer MB au dénominateur aurait associé les mauvais côtés.
\label{fig-cor-13-5}\figureGeo[corrige]{thales-05}
\end{corrige}\end{blocLivre}
\begin{blocLivre}{cor-13-6}\begin{corrige}{6}{Démontrer un parallélisme}\demarche{Vérifier les conditions géométriques avant les rapports.}
$AM/AB=4/10=2/5$ et $AN/AC=6/15=2/5$. Les alignements dans le même ordre sont donnés. La réciproque du théorème de Thalès donne $\res{(MN)\parallel(BC)}$.
\label{fig-cor-13-6}\figureGeo[corrige]{thales-06}
\end{corrige}\end{blocLivre}
\begin{blocLivre}{cor-13-7}\begin{corrige}{7}{Des droites non parallèles}\demarche{Comparer exactement, sans arrondir les quotients.}
$3/8\ne5/12$, car $3\times12=36$ et $5\times8=40$. D’après la contraposée du théorème de Thalès, les droites (MN) et (BC) $\res{\text{ne sont pas parallèles}}$.
\label{fig-cor-13-7}\figureGeo[corrige]{thales-07}
\end{corrige}\end{blocLivre}
\begin{blocLivre}{cor-13-8}\begin{corrige}{8}{Le papillon}\demarche{Conserver les longueurs positives dans les rapports.}
Thalès s’applique aussi au papillon : $AN/8=3/5$. Donc $\res{AN=4{,}8\text{ cm}}$. Les longueurs restent positives, même si les points sont de part et d’autre de A.
\label{fig-cor-13-8}\figureGeo[corrige]{thales-08}
\end{corrige}\end{blocLivre}
\begin{blocLivre}{cor-13-9}\begin{corrige}{9}{La distance entre deux points}\demarche{Lire l’ordre des points.}
A est entre les points de chaque paire : $\res{MB=3+5=8\text{ cm}}$ et $\res{NC=4{,}8+8=12{,}8\text{ cm}}$. Une soustraction correspondrait à une autre disposition.
\label{fig-cor-13-9}\figureGeo[corrige]{thales-09}
\end{corrige}\end{blocLivre}
\begin{blocLivre}{cor-13-10}\begin{corrige}{10}{Une réciproque en papillon}\demarche{Contrôler l’ordre des points, puis comparer les rapports.}
$AM/AB=2{,}4/4=0{,}6$ et $AN/AC=3/5=0{,}6$. Les points sont alignés dans le même ordre, de part et d’autre de A. La réciproque du théorème de Thalès donne $\res{(MN)\parallel(BC)}$.
\label{fig-cor-13-10}\figureGeo[corrige]{thales-10}
\end{corrige}\end{blocLivre}
\begin{blocLivre}{cor-13-11}\begin{corrige}{11}{Deux triangles de même aire}\demarche{Repérer une base commune et deux hauteurs égales.}
B et C sont sur une droite parallèle à (MN) : ils sont à la même distance de (MN). Les triangles MNB et MNC ont la même base [MN] et la même hauteur : $\res{\text{même aire}}$. Ensuite $\mathcal A(AMC)=\mathcal A(AMN)+\mathcal A(MNC)=\mathcal A(AMN)+\mathcal A(MNB)=\mathcal A(ANB)$. C’est l’idée de la preuve de Thalès par les aires.
\label{fig-cor-13-11}\figureGeo[corrige]{thales-aires}
\end{corrige}\end{blocLivre}
\begin{blocLivre}{cor-13-12}\begin{corrige}{12}{L’ombre d’un arbre}\demarche{Identifier le sommet commun et les deux parallèles.}
Le bâton et l’arbre sont verticaux, donc parallèles. Les pieds sont alignés avec S, ainsi que les sommets. Le théorème de Thalès dans les triangles emboîtés de sommet S donne $h/1{,}5=12/2$, d’où $\res{h=9\text{ m}}$. Les hypothèses de verticalité et de sol horizontal justifient le modèle.
\label{fig-cor-13-12}\figureGeo[corrige]{thales-ombre}
\end{corrige}\end{blocLivre}
\begin{blocLivre}{cor-13-13}\begin{corrige}{13}{Une équation dans Thalès}\demarche{Exprimer les côtés entiers avant de résoudre.}
AB = $4+x$ et AC = 9. Thalès donne $4/(4+x)=6/9=2/3$. Ainsi $12=8+2x$, donc $\res{x=2\text{ cm}}$. La valeur est positive et redonne $4/6=6/9$.
\label{fig-cor-13-13}\figureGeo[corrige]{thales-13}
\end{corrige}\end{blocLivre}
\begin{blocLivre}{cor-13-14}\begin{corrige}{14}{Calculer une largeur inaccessible}\demarche{Séparer la preuve du parallélisme et le calcul.}
M et N sont les milieux de deux côtés : par le théorème de la droite des milieux, ou par la réciproque de Thalès avec $AM/AB=AN/AC=1/2$, on a $(MN)\parallel(BC)$. Puis $MN/BC=1/2$, donc $\res{BC=14{,}8\text{ m}}$.
\label{fig-cor-13-14}\figureGeo[corrige]{thales-milieux}
\end{corrige}\end{blocLivre}
\begin{blocLivre}{cor-13-15}\begin{corrige}{15}{Un contrôle critique}\demarche{Vérifier que chaque terme est un côté correspondant.}
AM correspond à AN, donc $AN/16=3/8$, d’où AN = 6 cm. Puis $NC=16-6=\res{10\text{ cm}}$. NC n’est pas un côté du petit triangle AMN.
\label{fig-cor-13-15}\figureGeo[corrige]{thales-15}
\end{corrige}\end{blocLivre}
